Sampling in high dimension is a fundamental problem. Informally, given
access to a function $f:\R^{n}\rightarrow\R\cup\left\{ \infty\right\} $,
the sampling problem is to generate a point $x\in\R^{n}$ from the
distribution with density proportional to $e^{-f(x)}$. Note that
any density can be written in this form, so this is completely general.
To make the problem precise, we also have to specify a starting point
with positive density, and an error parameter $\epsilon$ that measures
the distance of the distribution of the output from the desired target. 

Unfortunately, in this generality, just like optimization, sampling
is also intractable. To see this, consider the following function:
\[
f(x)=\begin{cases}
0 & x\in S\\
M & x\not\in S
\end{cases}
\]
for some closed set $S$. Then sampling according to $e^{-f}$ for
a sufficiently large $M$ would allow us to find an element of $S$,
which could, e.g., be the minimizer of a hard-to-optimize function.

Consider a second example, which might appear more tractable: 
\[
g(x)=e^{-\frac{1}{2}x^{\top}Ax}\one_{x\ge0}.
\]
Without the restriction to the nonnegative orthant, the target density
is the Gaussian $N(0,A^{-1})$, and can be sampled by first sampling
the standard Gaussian $N(0,I)$ and applying the linear transformation
$A^{-1/2}$. To sample from the standard Gaussian in $\R^{n}$, we
can sample each coordinate independently from $N(0,1$), a problem
which has many (efficient) numerical recipes. But how can we handle
the restriction? In the course of forthcoming chapters, we will see
that this problem and its generalization to sampling logconcave densities,
i.e., when $f$ is convex, can be solved in polynomial time. It is
remarkable that the polynomial-time frontier for both optimization
and sampling is essentially determined by convexity.

We begin with gradient-based sampling methods. These rely on access
to $\nabla f$. These methods will in fact be natural algorithmic
versions of continuous processes on random variables, a particularly
pleasing connection. Later, we will see methods that only use access
to $f$, and others that utilize higher derivatives, notably the Hessian.
The parallels to optimization will be pervasive and striking.

\section{Gradient-based Sampling: Langevin Dynamics}

Here we study a simple stochastic process for generating samples from
a desired distribution $e^{-f(x)}$. As we will see later in this
chapter, it can also be viewed as a stochastic version of gradient
descent \emph{in the space of measures}. Gradient descent corresponds
to an Ordinary Differential Equation (ODE) for Gradient Flow:
\[
dX_{t}=-\nabla f(X_{t})\,dt
\]
which, by the chain rule, leads to 
\[
d(f(X_{t})-f(X_{0}))=\langle\nabla f(X_{t}),dX_{t}\rangle=-\norm{\nabla f(X_{t})}^{2}dt.
\]

\begin{xca}
Show that if $f$ is $\mu$-strongly convex, then under Gradient Flow,
we have 
\[
f(X_{t})-f(X^{*})\le e^{-2\mu t}(f(X_{0})-f(X^{*})).
\]
\end{xca}

For sampling from a given starting point, we need a process that introduces
randomness (since the target itself is a distribution). The continuous
version of such a process corresponds to a Stochastic Differential
Equation (SDE). It is called Langevin Diffusion and is the analog
of Gradient Flow for optimization.

\begin{algorithm2e}[H]

\caption{$\mathtt{LangevinDiffusion}$ (LD)}

\SetAlgoLined

\textbf{Input:} Initial point $x_{0}\in\Rn$.

Solve the stochastic differential equation
\[
dx_{t}=-\nabla f(x_{t})dt+\sqrt{2}dW_{t}
\]
\textbf{Output:} $x_{t}$.

\end{algorithm2e}

Here $f:\R^{n}\rightarrow\R$ is a differentiable function, $x_{t}$
is the random variable at time $t$ and $dW_{t}$ is infinitesimal
Brownian motion also known as a Wiener process. We can view it as
the limit of the following discrete process 
\[
x_{t+1}=x_{t}-h\nabla f(x_{t})+\sqrt{2h}\zeta_{t}
\]
with $\zeta_{t}$ sampled independently from $N(0,I)$. When we take
the step size $h\rightarrow0$, this discrete process converges to
the continuous one in distribution. We discuss the continuous version
first.

A more general form of an SDE is 
\begin{equation}
dx_{t}=\mu(x_{t},t)dt+\sigma(x_{t},t)dW_{t}\label{eq:sde}
\end{equation}
 where $x_{t}\in\R^{n},\mu(x_{t},t)\in\Rn$ is a time-varying vector
field and $\sigma(x_{t},t)\in\R^{n\times m}$ is a time-varying linear
transformation. The simplest such process is the Wiener process: $dx_{t}=dW_{t}$
which finds many applications in applied mathematics, finance, biology
and physics. Another useful process is the Ornstein-Ulhenbeck Process: 

\[
dx_{t}=-ax_{t}dt+\sigma dW_{t}
\]
This represents a particle moving in a fluid, the first term being
the force of friction exerted on the particle and the second term
being the movement caused by other particles colliding with our particle,
which is modeled using Brownian motion. 

A crucial difference between ordinary differentials and stochastic
differentials is in the chain rule. Unlike classical differentials
where we have $df(x)=\nabla f(x)dx$, we have the following chain
rule for stochastic calculus.
\begin{lem}[It\^{o}'s lemma]
\label{lem:itos}For any process $x_{t}\in\Rn$ satisfying $dx_{t}=\mu(x_{t})dt+\sigma(x_{t})dW_{t}$
where $\mu(x_{t})\in\Rn$ and $\sigma(x_{t})\in\R^{n\times m}$, we
have that
\begin{align*}
df(x_{t}) & =\nabla f(x_{t})^{\top}dx_{t}+\frac{1}{2}(dx_{t})^{\top}\nabla^{2}f(x_{t})(dx_{t})\\
 & =\nabla f(x_{t})^{\top}\mu(x_{t})dt+\nabla f(x_{t})^{\top}\sigma(x_{t})dW_{t}+\frac{1}{2}\tr(\sigma(x_{t})^{\top}\nabla^{2}f(x_{t})\sigma(x_{t}))dt.
\end{align*}
\end{lem}

Note that in Lemma \ref{lem:itos}, the dimension of the Brownian
motion need not match that of the process $x_{t}$ that it diffuses.
The usual chain rule comes from using Taylor expansion and taking
a limit, i.e., 
\[
\nabla f(x)=\lim_{h\rightarrow0}\frac{(f(x)+h\nabla f(x)^{\top}y+\frac{1}{2}h^{2}...)-f(x)}{h}=\lim_{h\rightarrow0}\frac{h\nabla f(x)^{\top}y+\frac{1}{2}h^{2}...}{h}
\]
and only the first term survives, as the second and later terms goes
to zero. But with the stochastic component, roughly speaking, $(dW_{t})^{2}=dt$\footnote{By definition, $W_{t}\sim N(0,t)$, so the variance of $W_{t}$ is
$t$. Thus, $(dW_{t})^{2}=dt$ may seem to make sense. This can be
justified rigorously through \emph{quadratic variations}.}, and we need to keep track of the second term in the Taylor expansion
as well. 

When $x_{t}$ is given by the SDE (\ref{eq:sde}), its \emph{quadratic
variation} $[x]_{t}$ is defined as
\[
[x]_{t}=\lim_{n\to\infty}\sum_{i=1}^{n}(x_{t_{i}}-x_{t_{i-1}})^{2}=\int_{0}^{t}\sigma_{s}^{2}\,ds,
\]
where $\{t_{i}\}_{0\leq i\leq n}$ is a partition of $[0,t]$ whose
mesh size $\max_{i\in[n]}|t_{i}-t_{i-1}|$ goes to zero as $n\to\infty$.
When $x_{t}$ is vector-valued, the quadratic variation is 
\[
[x]_{t}=\int_{0}^{t}\sigma_{s}\sigma_{s}^{\T}\,ds\,
\]
and $d[X^{i},X^{j}]_{t}=(\sigma_{t}\sigma_{t}^{\top})_{ij}dt$. We
can check that $[W]_{t}=tI$.

For a detailed treatment of stochastic calculus, we refer the reader
to a standard textbook such as \cite{oksendal2013stochastic}. 

First, we will see that $e^{-f}$ is a stationary density for the
Langevin SDE in continuous time. The proof relies on the following
general theorem about the distribution induced by an SDE.
\begin{thm}[Fokker--Planck equation]
\label{thm:fokker_planck}For any process $x_{t}\in\Rn$ satisfying
$dx_{t}=\mu(x_{t})dt+\sigma(x_{t})dW_{t}$ where $\mu(x_{t})\in\Rn$
and $\sigma(x_{t})\in\R^{n\times m}$ with the initial point $x_{0}$
drawn from $p_{0}$. Then the density $p_{t}$ of $x_{t}$ satisfies
the equation
\[
\frac{dp_{t}}{dt}=-\sum_{i}\frac{\partial}{\partial x_{i}}(\mu(x)_{i}p_{t}(x))+\frac{1}{2}\sum_{i,j}\frac{\partial^{2}}{\partial x_{i}\partial x_{j}}\left[(D(x))_{ij}p_{t}(x)\right]
\]
where $D(x)=\sigma(x)\sigma(x)^{\top}$.
\end{thm}

\begin{proof}
For any smooth function $\phi$, we have that
\[
\E_{x\sim p_{t}}\phi(x)=\E\phi(x_{t}).
\]
Taking derivatives on the both sides with respect to $t$, using It\^{o}'s
lemma (Lemma \ref{lem:itos}), and noting that $\E dW_{t}=0$, we
have that
\begin{align*}
\int\phi(x)\,dp(x)dx & =\E\left(\nabla\phi(x_{t})^{\top}\mu(x_{t})\,dt+\nabla\phi(x_{t})^{\top}\sigma(x_{t})\,dW_{t}+\frac{1}{2}\tr(\sigma(x_{t})^{\top}\nabla^{2}\phi(x_{t})\sigma(x_{t}))\,dt\right)\\
 & =\E\left(\nabla\phi(x_{t})^{\top}\mu(x_{t})dt+\frac{1}{2}\tr(\nabla^{2}\phi(x_{t})D(x_{t}))\,dt\right).
\end{align*}
Using $x_{t}\sim p_{t}$, we have that
\[
\int\phi(x)\frac{dp_{t}}{dt}dx=\int\nabla\phi(x)^{\top}\mu(x)p_{t}(x)\,dx+\frac{1}{2}\int\tr(\nabla^{2}\phi(x)D(x))p_{t}(x)\,dx.
\]
Integrating by parts, 
\[
\int\nabla\phi(x)^{\top}\mu(x)p_{t}(x)dx=-\int\phi(x)\sum_{i}\frac{\partial}{\partial x_{i}}(\mu_{i}(x)p_{t}(x))dx.
\]
Similarly, integrating by parts twice gives
\begin{align*}
\int\tr(\sigma(x)^{\top}\nabla^{2}\phi(x)\sigma(x))p_{t}(x)dx & =\int\tr(\nabla^{2}\phi(x)\sigma(x)\sigma(x)^{\top})p_{t}(x)dx\\
 & =-\int\langle\nabla\phi(x),\sum_{i}\frac{\partial}{\partial x_{i}}(p_{t}(x)D(x)_{i})\rangle dx\\
 & =\sum_{i,j}\int\phi(x)\frac{\partial^{2}}{\partial x_{i}\partial x_{j}}\left[(D(x))_{ij}p_{t}(x)\right]dx.
\end{align*}
Hence, 
\[
\int\phi(x)\left[\frac{dp_{t}}{dt}+\sum_{i}\frac{\partial}{\partial x_{i}}(\mu(x)_{i}p_{t}(x))-\frac{1}{2}\sum_{i,j}\frac{\partial^{2}}{\partial x_{i}\partial x_{j}}\left[(D(x))_{ij}p_{t}(x)\right]\right]dx=0
\]
for any smooth $\phi$. Therefore, we have the conclusion of the theorem.

We apply the Fokker--Planck equation to the Langevin dynamics.
\end{proof}
\begin{thm}
\label{thm:langevin_stationary}For any smooth function $f$, the
density proportional to $F=e^{-f}$ is stationary for the Langevin
dynamics.
\end{thm}

\begin{proof}
The Fokker--Planck equation (Theorem \ref{thm:fokker_planck}) shows
that the distribution $p_{t}$ of $x_{t}$ satisfies 
\begin{equation}
\frac{dp_{t}}{dt}=\sum_{i}\frac{\partial}{\partial x_{i}}(\frac{\partial f(x)}{\partial x_{i}}p_{t}(x))+\sum_{i}\frac{\partial^{2}}{\partial x_{i}^{2}}\left[p_{t}(x)\right].\label{eq:dp_LD}
\end{equation}
Now since $p_{t}$ is stationary the LHS is zero and we can rewrite
the above as
\begin{align*}
\frac{dp_{t}}{dt}=0 & =\sum_{i}\frac{\partial}{\partial x_{i}}\left(p_{t}(x)\frac{\partial f(x)}{\partial x_{i}}+\frac{\partial}{\partial x_{i}}p_{t}(x)\right)\\
 & =\sum_{i}\frac{\partial}{\partial x_{i}}\left(p_{t}(x)\left(\frac{\partial f(x)}{\partial x_{i}}+\frac{\partial}{\partial x_{i}}\log p_{t}(x)\right)\right)\\
 & =\sum_{i}\frac{\partial}{\partial x_{i}}\left(p_{t}(x)\left(\frac{\partial}{\partial x_{i}}\left[\log\frac{p_{t}(x)}{e^{-f(x)}}\right]\right)\right).
\end{align*}
We can verify that $p_{t}(x)\propto e^{-f(x)}$ is a solution.
\end{proof}
\begin{xca}
Consider the equation 
\[
\frac{d}{dt}p_{t}(x)=\frac{1}{2}\frac{d^{2}}{dx^{2}}p_{t}(x)
\]
 for $x\in\R,t>0$. Show that the density of $\mathcal{N}(0,t)$,
i.e. 
\[
\phi_{t}(x)=\frac{1}{\sqrt{2\pi t}}\exp\left(-\frac{x^{2}}{2t}\right)
\]
 satisfies this equation. Then generalize the process to $\R^{n}$.
\end{xca}

\begin{xca}
Consider the SDE with $X_{t}\in\R^{n}$: 
\[
dX_{t}=-X_{t}dt+\sqrt{2}dW_{t}.
\]
Use the Fokker-Planck equation to derive a corresponding stationary
density. Use It\^{o}'s lemma to derive $\E(\norm{X_{t}}_{4}^{4})$.
(Hint: Take expectation on both sides of It\^{o}'s lemma for appropriate
$f(X_{t})$, and use $\E df(X_{t})=d\E f(X_{t})$ for continuous $f$
and $df$). 
\end{xca}


\subsection{Convergence via Coupling in Wasserstein distance}

Next we turn to the rate of convergence, which will also prove uniqueness
of the stationary distribution for the stochastic process. For this,
we assume that $f$ is strongly convex. The proof is via the classical
coupling technique \cite{aldous1995reversible}. 

Our goal is to bound the rate at which the distribution of the current
point approaches the stationary distribution, in some chosen measure
of distance between distributions (for example, the TV distance).
To do this, in the coupling technique, we consider two points which
are both following the random process. One of them is already in the
stationary distribution, and therefore will stay there. The other
is our point. We will show that there is a coupling of the two distributions,
i.e., a joint distribution over the two points, whose marginals are
identical to the single point processes, such that the expected distance
between the two points decreases at a certain rate. More formally,
we couple two copies $x_{t},y_{t}$ of the random process with different
starting points (the coupling is a joint distribution $D(x_{t},y_{t})$
with the property that its marginal for each of $x_{t},y_{t}$ is
exactly the process) and show that their distributions get closer
over time. 

While the challenge usually is to find a good coupling, in the present
case, the simple identity coupling (i.e., the same Wiener process
is used for both $x_{t}$ and $y_{t}$) works well. The distance measure
we will use here is the Wasserstein distance (in Euclidean norm, see
Definition \ref{defn:The-Wasserstein--distance}). 
\begin{xca}
Show that for two distributions with the same finite support, computing
their Wasserstein distance reduces to a bipartite matching problem. 
\end{xca}

\begin{lem}
Let $x_{t},y_{t}$ evolve according to the Langevin diffusion for
a $\mu$-strongly convex function $f:\R^{n}\rightarrow\R$. Then,
there is a coupling $\gamma$ between $x_{t}$ and $y_{t}$ s.t.
\[
\E_{x_{t},y_{t}\sim\gamma}\norm{x_{t}-y_{t}}^{2}\le e^{-2\mu t}\norm{x_{0}-y_{0}}^{2}.
\]
\end{lem}

\begin{proof}
From the definition of LD, and by using the identity coupling, i.e.,
the same Gaussian $dW_{t}$ for both processes $x_{t}$ and $y_{t}$,
we have that 
\[
\frac{d}{dt}\left(x_{t}-y_{t}\right)=\nabla f(y_{t})-\nabla f(x_{t}).
\]
Hence, 
\[
\frac{1}{2}\frac{d}{dt}\|x_{t}-y_{t}\|^{2}=\left\langle \nabla f(y_{t})-\nabla f(x_{t}),x_{t}-y_{t}\right\rangle .
\]
Next, from the strong convexity of $f$, we have 
\begin{align*}
f(y_{t})-f(x_{t}) & \ge\nabla f(x_{t})^{\top}(y_{t}-x_{t})+\frac{\mu}{2}\norm{x_{t}-y_{t}}^{2},\\
f(x_{t})-f(y_{t}) & \ge\nabla f(y_{t})^{\top}(x_{t}-y_{t})+\frac{\mu}{2}\norm{x_{t}-y_{t}}^{2}.
\end{align*}
Adding two equations together, we have
\[
(\nabla f(x_{t})-\nabla f(y_{t}))^{\top}(x_{t}-y_{t})\geq\mu\norm{x_{t}-y_{t}}^{2}.
\]
Therefore, 
\[
\frac{1}{2}\frac{d}{dt}\|x_{t}-y_{t}\|^{2}\leq-\mu\|x_{t}-y_{t}\|^{2}.
\]
 Hence, 
\begin{align*}
\frac{d}{dt}\log\|x_{t}-y_{t}\|^{2} & =\frac{\frac{d}{dt}\|x_{t}-y_{t}\|^{2}}{\|x_{t}-y_{t}\|^{2}}\\
 & \leq-2\mu.
\end{align*}
Integrating both sides from $0$ to $t$, we get 
\[
\log\|x_{t}-y_{t}\|^{2}-\log\|x_{0}-y_{0}\|^{2}\le-2\mu t
\]
which proves the result.
\end{proof}
\begin{xca}
Give an example of a function $f$ for which the density proportional
to $e^{-f}$ is \emph{not} stationary for the following discretized
Langevin algorithm 
\[
x^{(k+1)}=x^{(k)}-\epsilon\nabla f(x^{(k)})+\sqrt{2\epsilon}Z^{(k)}
\]
where $Z^{(k)}\sim\mathcal{N}(0,1)$ are independent and the distribution
of $x^{(0)}$ is Gaussian.
\end{xca}


\subsection{Convergence in $\chi^{2}$ and KL divergences}

Next we analyze the convergence of Langevin dynamics in the two most
commonly used divergence metrics. The rate of mixing will depend on
a fundamental functional constants of the target density, which we
define next. These constants will play an important role in later
algorithms and analysis, and we will also see equivalent geometric
formulations of them.
\begin{defn}
We say that a probability measure $\nu$ on $\R^{n}$ satisfies a
\emph{Poincar\'e inequality} (PI) with parameter $C_{PI}(\nu)$ if
for all smooth functions $g:\R^{n}\to\R$,
\begin{equation}
\Var_{\nu}g\leq C_{PI}(\nu)\,\E_{\nu}\left(\norm{\nabla g}^{2}\right)\label{eq:pi}
\end{equation}
where $\Var_{\nu}g=\E_{\nu}\abs{g-\E_{\nu}g}^{2}$.
\end{defn}

The Poincar\'e inequality is implied by the generally stronger log-Sobolev
inequality.
\begin{defn}
We say that a probability measure $\nu$ on $\R^{n}$ satisfies a
\emph{log-Sobolev inequality} (LSI) with parameter $C_{LSI}(\nu)$
if for all smooth functions $g:\Rn\to\R$,
\begin{equation}
\Ent_{\nu}(g^{2})\leq2C_{LSI}(\nu)\,\E_{\nu}\left(\norm{\nabla g}^{2}\right)\label{eq:lsi}
\end{equation}
where $\Ent_{\nu}(g^{2}):=\E_{\nu}(g^{2}\log g^{2})-\E_{\nu}(g^{2})\log\E_{\nu}(g^{2})$. 
\end{defn}

\begin{thm}
\label{thm:LDinKL}Let $f$ be a smooth function with log-Sobolev
constant $C_{LSI}$. Then the Langevin dynamics 
\[
dx_{t}=-\nabla f(x)dt+\sqrt{2}dW_{t}
\]
converges exponentially in KL-divergence (Definition~\ref{defn:KL-divergence})
to the density $\nu(x)\propto e^{-f(x)}$ with mixing rate $O(C_{LSI})$,
i.e., the distribution $\rho_{t}$of $x_{t}$satisfies 
\[
d_{KL}(\rho_{t},\nu)\le e^{-\frac{2}{C_{LSI}}t}d_{KL}(\rho_{0},\nu).
\]
\end{thm}

The theorem will follow from the following lemma and the log-Sobolev
inequality.
\begin{lem}
For a target distribution $\nu$, with $\rho_{t}$ being the distribution
attained by Langevin dynamics at time $t$, starting at $\rho_{0}$,
we have
\[
\frac{d}{dt}d_{KL}(\rho_{t},\nu)=-J_{\nu}(\rho_{t})
\]
where $J_{\nu}(\rho)$ is the relative Fisher information of $\rho$
wrt $\nu$, defined as $\E_{\rho}(\norm{\nabla\log(\rho/\nu)}^{2})$.
\end{lem}

\begin{proof}
We compute the time derivative of the distribution of LD at time $t$
with $\rho=\rho_{t}$:
\begin{align*}
\frac{d}{dt}d_{KL}(\rho_{t},\nu) & =\frac{d}{dt}\int_{\R^{n}}\rho\log\frac{\rho}{\nu}\\
 & =\int\frac{d\rho}{dt}\left(\log\frac{\rho}{\nu}+1\right)\\
 & =\int\nabla\cdot\left(\rho\nabla\log\frac{\rho}{\nu}\right)\log\frac{\rho}{\nu}\\
 & =-\int\langle\rho\nabla\log\frac{\rho}{\nu},\nabla\log\frac{\rho}{\nu}\rangle\\
 & =-\int\rho\norm{\nabla\log\frac{\rho}{\nu}}_{2}^{2}\\
 & =\E_{\rho}\left(\norm{\nabla\log\frac{\rho}{\nu}}_{2}^{2}\right)=-J_{\nu}(\rho_{t}).
\end{align*}
\end{proof}
The next lemma shows that the log-Sobolev constant can be equivalently
formulated in terms of distributions in place of smooth functions.
\begin{lem}
The log-Sobolev constant of a distribution $\nu$ on $\R^{n}$can
be equivalently defined as follows: for any distribution $\rho$ on
$\R^{n}$, we have
\[
d_{KL}(\rho,\nu)\le\frac{C_{LSI}}{2}J_{\nu}(\rho).
\]
\end{lem}

\begin{proof}
To go one way, we use $g=\sqrt{\frac{\rho}{\nu}};$to go the other
way, define $\rho=\frac{g^{2}\nu}{\E_{\nu}(g^{2})}$.
\end{proof}
We are now ready to prove Theorem \ref{thm:LDinKL}.
\begin{proof}
We use Lemmas and to obtain
\[
\frac{d}{dt}d_{KL}(\rho_{t},\nu)=-J_{\nu}(\rho_{t})\le-\frac{2}{C_{LSI}}d_{KL}(\rho_{t},\nu).
\]
Integrating over time gives the theorem.
\end{proof}
\begin{xca}
Show that for Langevin Dynamics applied to a target density $\nu$
with Poincar\'e constant $C_{PI}$ satisfies 
\[
\chi^{2}(\rho_{t},\nu)\le e^{-\frac{2}{C_{PI}}t}\chi^{2}(\rho_{0},\nu).
\]
\end{xca}

\section{Langevin Dynamics is Gradient Descent in Density Space*\footnote{Sections marked with * are more mathematical and can be skipped.}}\label{sec:LD_GD}

Here we show that Langevin dynamics is simply gradient descent for
the function $F(\rho)=\KL(\rho\|\nu)$ for two densities $\rho,\nu$
on the Wasserstein space where $\nu=e^{-f(x)}/\int e^{-f(y)}dy$ is
the target and $\rho$ is the current density. For this, we first
define the Wasserstein space.
\begin{defn}
The Wasserstein space $P_{2}(\Rn)$ on $\Rn$ is the manifold on the
set of probability measures on $\Rn$ such that the shortest path
distance of two measures $x,y$ in this manifold is exactly equal
to the Wasserstein distance between $x$ and $y$.

We let $T_{p}({\cal M})$ refer to the tangent space at a point $p$
in a manifold ${\cal M}.$
\end{defn}

\begin{lem}
\label{lem:local-norm}For any $p\in P_{2}(\Rn)$ and $v\in T_{p}P_{2}(\Rn)$,
we can write $v(x)=\nabla\cdot(p(x)\nabla\lambda(x))$ for some function
$\lambda$ on $\Rn$. Furthermore, the local norm of $v$ in this
metric is given by
\[
\|v\|_{p}^{2}=\E_{x\sim p}\|\nabla\lambda(x)\|^{2}.
\]
\end{lem}

\begin{proof}
Let $p\in P_{2}(\Rn)$ and $v\in T_{p}P_{2}(\Rn)$. We will show that
any change of density $v$ can be represented by a vector field $c$
on $\Rn$ as follows: Consider the process $x_{0}\sim p$ and $\frac{d}{dt}x_{t}=c(x_{t})$.
Let $p_{t}$ be the density of the distribution of $x_{t}$. To compute
$\frac{d}{dt}p_{t}$, we follow the same idea as in the proof as Theorem
\ref{thm:fokker_planck}. For any smooth function $\phi$, we have
that $\E_{x\sim p_{t}}\phi(x)=\E\phi(x_{t}).$ Taking derivatives
on the both sides with respect to $t$, we have that
\[
\int\phi(x)\frac{d}{dt}p_{t}(x)dx=\int\nabla\phi(x)^{\top}c(x)p_{t}(x)dx=-\int\nabla\cdot(c(x)p_{t}(x))\phi(x)dx
\]
where we used integration by parts at the end. Since this holds for
all $\phi$, we have that
\[
\frac{dp_{t}(x)}{dt}=-\nabla\cdot(p_{t}(x)c(x)).
\]
Since we are interested only in vector fields that generate the minimum
movement in Wasserstein distance, we consider the optimization problem
\[
\min_{-\nabla\cdot(pc)=v}\frac{1}{2}\int p(x)\|c(x)\|^{2}dx
\]
where we can think $v$ is the change of $p_{t}$. Let $\lambda(x)$
be the Lagrangian multiplier of the constraint $-\nabla\cdot(pc)=v$.
Then, the problem becomes
\begin{align*}
 & \min_{c}\frac{1}{2}\int p(x)\|c(x)\|^{2}dx-\int\lambda(x)\nabla\cdot(p(x)c(x))dx.\\
= & \min_{c}\frac{1}{2}\int p(x)\|c(x)\|^{2}dx+\int\nabla\lambda(x)^{\top}c(x)\cdot p(x)dx.
\end{align*}
Now, we note that the problem is a pointwise optimization problem
whose minimizer is given by
\[
c(x)=-\nabla\lambda(x).
\]
This proves that any vector field that generates minimum movement
in Wasserstein distance is a gradient field. Also, we have that $v(x)=\nabla\cdot(p(x)\nabla\lambda(x)).$
Note that the right hand side is an elliptical differential equation
and hence for any $v$ with $\int v(x)dx=0$, there is an unique solution
$\lambda(x)$. Therefore, we can write $v(x)=\nabla\cdot(p(x)\nabla\lambda(x))$
for some $\lambda(x)$.

Next, we note that the movement is given by
\[
\|v\|_{p}^{2}=\int p(x)\|c(x)\|^{2}dx=\E_{x\sim p}\|\nabla\lambda(x)\|^{2}.
\]
\end{proof}
As we discussed in the gradient descent section, one can use norms
other than $\ell_{2}$ norm. For the Wasserstein space, we should
use the local norm as given in Lemma \ref{lem:local-norm}.
\begin{thm}
\label{thm:LD_GD}Let $\rho_{t}$ be the density of the distribution
produced by Langevin Dynamics for the target distribution $\nu=e^{-f(x)}/\int e^{-f(y)}dy$.
Then, we have that
\[
\frac{d\rho}{dt}=\argmin_{v\in T_{p}P_{2}(\Rn)}\left\langle \nabla F(\rho),v\right\rangle _{p}+\frac{1}{2}\|v\|_{p}^{2}.
\]
Namely, $\rho_{t}$ follows continuous gradient descent in the density
space for the function $F(\rho)=\KL(\rho\|\nu)$ under the Wasserstein
metric.
\end{thm}

\begin{proof}
For any function $c$, the optimization problem of interest satisfies
\[
\min_{\delta=\nabla\cdot(\rho\nabla\lambda)}\left\langle c,\delta\right\rangle +\frac{1}{2}\int\rho(x)\|\nabla\lambda(x)\|^{2}dx=\min_{\nabla\lambda}-\int\rho(x)\cdot\nabla c(x)^{\top}\nabla\lambda(x)dx+\frac{1}{2}\int\rho(x)\|\nabla\lambda(x)\|^{2}dx.
\]
Solving the right hand side, we have $\nabla c=\nabla\lambda$ and
hence $\delta=\nabla\cdot(\rho\nabla c)$. Now, we note that $\nabla F(\rho)=\log\frac{\rho}{\nu}-1$.
Therefore, 
\begin{align*}
\frac{d\rho}{dt} & =\nabla\cdot(\rho\nabla(\log\frac{\rho}{\nu}-1))\\
 & =\nabla\cdot(\rho\nabla\log\frac{\rho}{\nu})\\
 & =\nabla\cdot\left(\rho\nabla f\right)+\Delta\rho
\end{align*}
which is exactly equal to (\ref{eq:dp_LD}). 
\end{proof}
To analyze this continuous descent in Wasserstein space, we first
prove that continuous gradient descent converges exponentially whenever
$F$ is strongly convex.
\begin{lem}
\label{lem:grad-dom}Let $F$ be a function satisfying ``Gradient
Dominance'':
\begin{equation}
\norm{\nabla F(x)}_{x}^{2}\ge\alpha\cdot(F(x)-\min_{y}F(y))\quad\text{for all }x\label{eq:LD_strong_convexity}
\end{equation}
on the manifold with the metric $\|\cdot\|_{x}$ where $\nabla$ is
the gradient on the manifold. Then, the process $dx_{t}=-\nabla F(x_{t})dt$
converges exponentially, i.e., $F(x_{t})-\min_{y}F(y)\le e^{-\alpha t}(F(x_{0})-\min_{y}F(y))$.
\end{lem}

\begin{proof}
We write
\[
\frac{d}{dt}(F(x)-\min_{y}F(y))=\langle\nabla F(x_{t}),\frac{dx_{t}}{dt}\rangle_{x_{t}}=-\|\nabla F(x_{t})\|_{x_{t}}^{2}\leq-\alpha(F(x)-\min_{y}F(y)).
\]
The conclusion follows.
\end{proof}
Finally, we note that the log-Sobolev inequality for the density $\nu$
can be re-stated as the condition (\ref{eq:LD_strong_convexity}).
\begin{lem}
\label{lem:log_sob_LD}Fix a density $\nu$. Then the log-Sobolev
inequality, namely, for every smooth function $g$, 
\[
\frac{C_{LSI}}{2}\int\norm{\nabla g}^{2}\,d\nu\ge\int g(x)^{2}\log g(x)^{2}\ d\nu
\]
 implies the condition (\ref{eq:LD_strong_convexity}) with $\alpha=1/C_{LSI}$.
\end{lem}

\begin{proof}
Take $g(x)=\sqrt{\frac{\rho(x)}{\nu(x)}}$, the log-Sobolev inequality
shows that 
\[
\frac{C_{LSI}}{2}\int\rho(x)\norm{\nabla\log\frac{\rho(x)}{\nu(x)}}^{2}\,dx\geq\int\rho(x)\log\frac{\rho(x)}{\nu(x)}\,dx\text{ for all }\rho.
\]

As we calculate in Theorem \ref{thm:LD_GD}, we have that
\[
\norm{\nabla F(\rho)}_{\rho}^{2}=\int\rho(x)\norm{\nabla\log\frac{\rho(x)}{\nu(x)}}^{2}\,dx.
\]
Therefore, this is exactly the condition (\ref{eq:LD_strong_convexity})
with coefficient $2/C_{LSI}$.
\end{proof}
Combining Lemma \ref{lem:log_sob_LD} and Lemma \ref{lem:grad-dom},
we obtain another proof of Therem \ref{thm:LDinKL}.

\subsection{Discussion}

Langevin dynamics converges quickly in continuous time for isoperimetric
distributions. Turning this into an efficient algorithm typically
needs more assumptions and there is much room for choosing discretizations.
This is similar to the situation with gradient descent for optimization.
As we saw in Section \ref{sec:LD_GD}, it turns out that Langevin
dynamics is in fact gradient descent in the space of probability measures
under the Wasserstein metric, where the function being minimized is
the KL-divergence of the current density from the target stationary
density. For more on this view of sampling as optimization over measures,
see \cite{wibisono2018sampling}. See \cite{lee2024eldan} for a tight
estimate of log-Sobolev constant for logconcave measures. In particular
for a logconcave measure with support of diameter $D$, the log-Sobolev
constant is $\Omega(1/D)$.
